\documentclass[aspectratio=169,10pt]{beamer}
\usepackage[utf8]{inputenc}
\usepackage[T1]{fontenc}
\usepackage[english]{babel}
\babelprovide[import,main]{french}
\usepackage{lmodern,amsmath,amssymb,booktabs,graphicx}
\newcommand{\DataCommit}{57ab54ebc994029be3c24a8a34da56371547f4ac}
\newcommand{\DataURL}[1]{https://github.com/halhoulmhamed-droid/inverse-laplace-learning/blob/57ab54ebc994029be3c24a8a34da56371547f4ac/studies/study2/#1}

\definecolor{ink}{RGB}{25,48,73}
\definecolor{accent}{RGB}{23,116,126}
\setbeamercolor{normal text}{fg=ink,bg=white}
\setbeamercolor{frametitle}{fg=ink}
\setbeamercolor{structure}{fg=accent}
\setbeamercolor{block title}{bg=accent!12,fg=ink}
\setbeamercolor{block body}{bg=accent!4,fg=ink}
\setbeamertemplate{navigation symbols}{}
\setbeamertemplate{footline}{\hfill\scriptsize Halhoul Mhamed\quad Étude 2\quad\insertframenumber/\inserttotalframenumber\hspace{5mm}\vspace{3mm}}
\setbeamersize{text margin left=9mm,text margin right=9mm}
\hypersetup{pdftitle={Étude 2 : covariances apprises et observations partielles pour Laplace},pdfauthor={Halhoul Mhamed}}
\newcommand{\src}[2]{\href{\DataURL{#1}}{#2}}
\newcommand{\figsource}[1]{\par\tiny\src{figures/#1.npz}{Données}\,|\,\src{src/figures.py}{Script}\,|\,\src{results/final/configuration.json}{Configuration}\,|\,\src{figures/#1.pdf}{Version scientifique}\par}
\newcommand{\plot}[2]{\includegraphics[width=#2\linewidth]{../figures/#1.pdf}}
\begin{document}
\begin{frame}[plain]
\begin{columns}[T]
\column{.48\textwidth}\includegraphics[width=\linewidth,height=12mm,keepaspectratio]{../../../assets/logos/university.png}
\column{.48\textwidth}\raggedleft\includegraphics[width=\linewidth,height=12mm,keepaspectratio]{../../../assets/logos/faculty.png}
\end{columns}
\centering
{\scriptsize Université Abdelmalek Essaâdi --- Faculté Polydisciplinaire de Larache\par}
\par\vspace{2mm}
{\large Projet scientifique de Master --- Étude 2\par}
\par\vspace{3mm}
{\LARGE\bfseries Complétion de données de Laplace\par}
{\large Apport des covariances apprises et robustesse\par aux observations partielles\par}
\par\vspace{4mm}
{\large Halhoul Mhamed\par}
{\scriptsize Master « Modélisation mathématique et applications et apprentissages » --- M2 en cours\par}
\par\vspace{4mm}
\textbf{Professeur encadrant : Pr. Chakir Tajani}\par
\vfill
{\tiny 11 octobre 2026}
\end{frame}

\begin{frame}{Question et réponse limitée}
\begin{block}{Question préspécifiée}
Les covariances complètes améliorent-elles les covariances diagonales dans les mêmes 12 directions singulières, avec mêmes données, moyennes et réglage ?
\end{block}
\par\vspace{3mm}
\begin{itemize}
\item Oui selon la règle figée : risque F/D proche de 0,32 sur 32 capteurs contigus, $N=256$.
\item LMMSE complet bénéficie aussi de la prédiction hors du sous-espace : effet distinct.
\item Dégradations observées face à H1 avec tous les capteurs et sur des objets plus rugueux.
\end{itemize}
\par\vspace{2mm}
\small Les résultats de l'étude 1 ont formulé l'hypothèse ; ses objets de test ne sont pas réutilisés. Sa conclusion négative reste intégralement conservée.
\end{frame}

\begin{frame}{Modèle : indépendance dans la base cosinus}
\[
q(x)=\sum_{k=0}^{127}c_k\phi_k(x),\quad c_k=w_k\xi_k,\quad
\xi_k\ \text{indépendants},\quad w_k^2\propto(1+k^2)^{-p}.
\]
\begin{itemize}
\item Rectangle harmonique, $H=0,25$, 64 moyennes de capteurs ; transfert $\operatorname{sech}(k\pi H)$.
\item Générateur analytique 128 modes ; inverse FD à 64 modes, grille $64\times16$ intervalles.
\item Bruit absolu gaussien ; observations divisées par 8, poids $1/64$ conservés après masque.
\end{itemize}
\begin{block}{Attention à l'interprétation}
$\operatorname{Cov}(V_r^\top c)=V_r^\top D_cV_r$ peut être non diagonale. Changement de base et capteurs partiels ne prouvent pas des corrélations physiques des coefficients cosinus.
\end{block}
\end{frame}

\begin{frame}{Cinq restrictions pour séparer les effets}
\small $z=U_r^\top(y-\mu_y)$, $b=V_r^\top(c-\mu_c)$, $r=12$, $P=V_rV_r^\top$. Toujours $\widehat c=\mu_c+R(y-\mu_y)$.
\par\vspace{3mm}
\begin{center}
\begin{tabular}{llll}
\toprule & Entrée & Sortie de l'incrément & Covariances\\\midrule
D & $U_r$ & $V_r$ & Diagonales $C_{bz},C_{zz}$\\
F & $U_r$ & $V_r$ & Blocs complets\\
O & $U_r$ & Toute la sortie & Blocs complets\\
L & Toutes les observations & Toute la sortie & Blocs complets\\
LP & Toutes les observations & $V_r$ & Projection de L\\\bottomrule
\end{tabular}
\end{center}
\par\vspace{3mm}
\begin{block}{Identités testées à même ridge}
$PR_O=R_F$ et $PR_L=R_{LP}$. Projeter l'incrément, pas la moyenne. Les diagnostics utilisent $\lambda_F$ pour D/F/O et $\lambda_L$ pour L/LP.
\end{block}
\end{frame}

\begin{frame}{Pourquoi l'ablation est mathématiquement contrôlée}
\[
\min_G\;\mathbb E\|b-Gz\|^2+\lambda t\|G\|_F^2
\quad\Longrightarrow\quad G=C_{bz}(C_{zz}+\lambda tI)^{-1}.
\]
\begin{itemize}
\item Contrainte $G$ diagonale : $g_j=C_{bz,jj}/(C_{zz,jj}+\lambda t)$.
\item Même base, même sous-espace, même échelle $t$, mêmes moyennes, 20 candidats chacun.
\item Version empirique régularisée de LMMSE ; aucune optimalité universelle sur le test.
\end{itemize}
\[
\|\widehat q-q\|^2=\|Pe\|^2+\|(I-P)e\|^2+\sum_{k=64}^{127}c_k^2.
\]
\small À ridge commun, F et O ont exactement la même perte dans $P$. Leur écart total provient du complément, qui n'est pas nécessairement le noyau exact de l'opérateur.
\end{frame}

\begin{frame}{Protocole figé avant l'évaluation finale}
\begin{columns}[T]
\column{.51\textwidth}
\textbf{Données nouvelles}
\begin{itemize}
\item $N=64,256,1024$ ; cinq entraînements indépendants.
\item 256 objets validation, deux bruits ; réglage supervisé uniquement sur validation.
\item 400 objets test, quatre bruits appariés. Partitions pilote distinctes.
\end{itemize}
\column{.47\textwidth}
\textbf{Décision principale}
\begin{itemize}
\item F/D, 32 contigus, $N=256$, bruit 0,01, nominal.
\item Risque quadratique : ratio $\leq0,95$.
\item Protection $q_{95}$ : ratio $\leq1,05$.
\item 4\,000 bootstraps par objet ; bornes 99,5\,\%, deux critères dans cinq jeux.
\end{itemize}
\end{columns}
\par\vspace{3mm}
\small Gel : \texttt{20a663d}. Couverture simultanée nominale approximative ; incertitude conditionnelle aux modèles. Pas d'intervalle populationnel sur les entraînements.
\end{frame}

\begin{frame}{Contraste principal : critères satisfaits}
\centering\plot{primary}{.90}\figsource{primary}
\par\vspace{3mm}
\normalsize Les cinq bornes de risque sont $<0,336$ ; les cinq bornes de queue sont $<0,375$. Les seuils 0,95 et 1,05 sont respectés.
\par\vspace{2mm}
\small Conclusion restreinte au scénario et aux cinq modèles observés ; réduction d'erreur quadratique, pas garantie individuelle.
\end{frame}

\begin{frame}{Huit méthodes : ne pas confondre F et LMMSE complet}
\begin{columns}[c]
\column{.54\textwidth}
\centering\small\begin{tabular}{lrr}
\toprule Méthode & Risque moyen & $q_{0,95}$ moyen\\\midrule
D & 0.160046 & 0.583969\\
F & 0.051333 & 0.180578\\
O & 0.012028 & 0.037325\\
L & 0.012388 & 0.038519\\
LP & 0.051577 & 0.181273\\
H1 & 0.016452 & 0.055207\\
Signé plafonné & 0.160076 & 0.581725\\
Positif borné & 0.159963 & 0.582030\\
\bottomrule
\end{tabular}

\column{.43\textwidth}
\small Scénario principal, $N=256$. Moyennes sur cinq entraînements ; chaque méthode est réglée séparément.
\par\vspace{4mm}
\normalsize F améliore D mais reste derrière H1.
\par\vspace{3mm}
L/H1 : ratio observé 0,753. O fait légèrement mieux que L ici.
\par\vspace{3mm}
\small La comparaison L/H1 est secondaire et descriptive. La colonne $q_{95}$ est la moyenne des cinq percentiles.
\end{columns}
\end{frame}

\begin{frame}{Tailles d'apprentissage : dispersion distincte de l'incertitude}
\centering\plot{training}{1}\figsource{training}
\par\vspace{3mm}
\small Cinq courbes d'entraînement, tailles imbriquées. Les objets de test restent communs ; les courbes ne sont pas des intervalles de confiance populationnels.
\end{frame}

\begin{frame}{Une contribution séparée : la sortie hors sous-espace}
\centering\plot{projection}{.96}\figsource{projection}
\par\vspace{2mm}
\small À ridge partagé, F/O puis L/LP isolent la sortie complémentaire. Les moyennes sont identiques. Le complément spectral n'est pas assimilé au seul noyau de $A$.
\end{frame}

\begin{frame}{Robustesse : les dégradations restent visibles}
\begin{columns}[c]
\column{.72\textwidth}
\plot{robustness}{1}\figsource{robustness}
\column{.26\textwidth}
\small L face à H1 :
\par\vspace{2mm}
\textbf{+9,2\,\%} de risque avec tous les capteurs.
\par\vspace{3mm}
\textbf{+10,7\,\%} sur objets rugueux.
\par\vspace{3mm}
16 contigus : risque moyen réduit, mais $q_{95}$ moyen \textbf{+4,4\,\%}.
\par\vspace{3mm}
Ratios descriptifs ; pas de preuve corrigée de supériorité secondaire.
\end{columns}
\end{frame}

\begin{frame}{Covariance de population et effet de coordonnées}
\centering\plot{covariance}{1}\figsource{covariance}
\par\vspace{3mm}
\small $V_r^\top D_cV_r$ est calculée à partir des variances connues, sans estimation sur les vérités de test. « Hors diagonale » désigne une fraction de norme de Frobenius. Les coefficients cosinus restent indépendants.
\end{frame}

\begin{frame}{Contrôles réussis et limites de portée}
\begin{columns}[T]
\column{.49\textwidth}
\textbf{Vérifications}
\begin{itemize}
\item 11 tests ciblés ; pilote séparé réussi.
\item 3,12 millions d'évaluations ; aucune panne numérique.
\item 468 valeurs recomposées indépendamment, écart $<2\times10^{-15}$.
\item 209 fichiers de l'étude 1 inchangés octet pour octet.
\end{itemize}
\column{.49\textwidth}
\textbf{Limites}
\begin{itemize}
\item Un seul rang $r=12$ ; données synthétiques gaussiennes.
\item Validation supervisée et cinq entraînements seulement.
\item Grilles finies, sélections aux bornes conservées.
\item Pas de bruit corrélé, de valeurs aberrantes ni de géométrie différente.
\end{itemize}
\end{columns}
\par\vspace{3mm}
\small Campagne seule : 6,83 s CPU, pic 83,1 Mo, sorties 67,5 Mo. Plafonds globaux : 4 h CPU, 2 Gio RAM, 200 Mo de nouvelles sorties. Aucun service payant, GPU ou FreeFem.
\end{frame}

\begin{frame}{Conclusion, antériorité et contact}
\begin{block}{Conclusion}
L'apport des covariances complètes est établi pour le contraste figé à espace identique. La sortie hors sous-espace apporte un gain distinct. Aucune supériorité générale de LMMSE n'est établie.
\end{block}
\small
Antériorités : Burger--Kabri (2024), Alberti et al. (NeurIPS 2021), Hansen (Regularization Tools), Tajani--Abouchabaka (2018). \src{SOURCES.md}{Références et limites de lecture}. Méthodes établies, aucune nouveauté algorithmique revendiquée.
\par\vspace{2mm}
Code, analyse et rédaction avec assistance d'IA ; validation académique relevant de l'auteur et de l'encadrant. Les logos ne valent pas approbation des résultats.
\par\vspace{3mm}
\textbf{Halhoul Mhamed}\par
\href{mailto:halhoul.mhamed@etu.uae.ac.ma}{halhoul.mhamed@etu.uae.ac.ma}\par
\href{mailto:halhoulmhamed@gmail.com}{halhoulmhamed@gmail.com}\par
\href{https://github.com/halhoulmhamed-droid/inverse-laplace-learning}{Dépôt public inverse-laplace-learning}\par
\par\vspace{2mm}
{\tiny Version scientifique : \texttt{\DataCommit}}
\end{frame}
\end{document}
